% ----------------------------------------------------------------------------
\chapter{The sharp isoperimetric inequality}\label{ch:isoperimetric}
\module{NoCompromise.Isoperimetric.ABP}

\begin{remark}[Acyclicity]\label{rem:acyclic-isoperimetric}
The \emph{inequality} \eqref{eq:sharp-isoperimetric} is proved here from the
ABP contact-set argument and depends on nothing from
Part~\ref{part:regularity}.  The \emph{equality case} needs the regularity
theory for $\omega$-minimal sets and is therefore deferred to
Theorem~\ref{thm:isoperimetric-rigidity}.  Chapter~\ref{ch:penalization}
onwards uses only the inequality.  This ordering is what keeps the dependency
graph acyclic, and it must be preserved: if the regularity chapters are ever
allowed to cite the equality case, the development becomes circular.
\end{remark}

\begin{lemma}[The ABP auxiliary Neumann problem]\label{lem:abp-neumann}
\uses{prop:weak-neumann,thm:boundary-neumann}
Let $G\subset\R^3$ be a bounded connected smooth domain.  There is a smooth
$z:\overline G\to\R$ with
\begin{equation}\label{eq:neumann-abp}
  \Delta z=\frac{\Per(G)}{|G|}\ \text{in }G,
  \qquad
  \partial_\nu z=1\ \text{on }\partial G .
\end{equation}
\end{lemma}

\begin{proof}\uses{prop:weak-neumann,thm:boundary-neumann}
Compatibility \eqref{eq:neumann-compat} holds because
$\int_G\Per(G)/|G|=\Per(G)=\int_{\partial G}1$.  Apply
Proposition~\ref{prop:weak-neumann} with $f=\Per(G)/|G|$, $h=1$ and the sign
convention $\partial_\nu z=h$.  In an upper-half-ball signed-distance chart,
the inward-coordinate conormal appearing in
Theorem~\ref{thm:boundary-neumann} is the negative of this outward datum
(with the smooth chart density absorbed into $A,f,h$).  To check ellipticity,
write the chart as $x=\Psi(y)$, set $J=|\det D\Psi|$, and use the pulled-back
Laplacian coefficient
\[
 A=J(D\Psi)^{-1}(D\Psi)^{-\mathsf T},
 \qquad \xi\cdot A\xi=J|(D\Psi)^{-\mathsf T}\xi|^2.
\]
On finitely many smaller chart closures, $J$ is bounded above and below by
positive constants, and $D\Psi$ and its inverse are bounded.  This gives
uniform $0<\lambda\le\Lambda<\infty$.  Signed-distance coordinates make the
normal direction orthogonal to the tangential ones on the flat face, giving
$A_{3\alpha}=A_{\alpha3}=0$ there.  The smooth boundary gives coefficients
and data of every required order, so Theorem~\ref{thm:boundary-neumann}
applies and gives smoothness up to the boundary.
\end{proof}

\begin{lemma}[The contact set covers the unit ball of gradients]
\label{lem:abp-contact}
\uses{lem:abp-neumann}
Let $z$ be as in Lemma~\ref{lem:abp-neumann} and let
\[
  \Gamma:=\{x\in G: z(y)\ge z(x)+\nabla z(x)\cdot(y-x)\ \ \forall y\in\overline G\}
\]
be its lower contact set.  Then $B_1\subset\nabla z(\Gamma)$, and
$D^2z\ge0$ on $\Gamma$.
\end{lemma}

\begin{proof}\uses{lem:abp-neumann}
For $p\in B_1$ the function $z-p\cdot x$ attains its minimum on
$\overline G$.  A boundary minimum would force
$\partial_\nu(z-p\cdot x)\le0$, but
$\partial_\nu(z-p\cdot x)=1-p\cdot\nu>0$.  So the minimum is interior, where
$\nabla z=p$ and $D^2z\ge0$; such a point lies in $\Gamma$ by convexity of the
supporting-plane condition at an interior minimum.
\end{proof}

\begin{proposition}[Smooth connected case]\label{prop:iso-smooth}
\uses{lem:abp-contact}
For a bounded connected smooth domain $G\subset\R^3$,
\begin{equation}\label{eq:iso-smooth}
  \Per(G)^3\ge36\pi\,|G|^2 .
\end{equation}
\end{proposition}

\begin{proof}\uses{lem:abp-contact,lem:abp-neumann,lem:sard-equidim}
The critical values of $\nabla z$ are null by
Lemma~\ref{lem:sard-equidim}.  Cover the regular set by countably many inverse
function neighbourhoods and disjointify them; on each resulting Borel piece
the equidimensional change-of-variables theorem from \ref{base:calculus}
applies injectively.  Summing gives
$|\nabla z(\Gamma)|\le\int_\Gamma|\det D^2z|$.  On $\Gamma$ the Hessian is
positive semidefinite, so Lemma~\ref{lem:abp-contact} and the
arithmetic--geometric mean inequality give
\[
  |B_1|\le\int_\Gamma\det D^2z
  \le\int_\Gamma\Bigl(\frac{\Delta z}{3}\Bigr)^3
  \le|G|\Bigl(\frac{\Per(G)}{3|G|}\Bigr)^3 .
\]
Since $|B_1|=4\pi/3$, this is \eqref{eq:iso-smooth}.
\end{proof}

\medskip\noindent\emph{Formalisation note.} The argument uses only that $\partial G$ is $C^3$: the Neumann solution $z$ of Lemma~\ref{lem:abp-neumann} is then $C^2$ in $G$ and $C^1$ up to the boundary with $\partial_\nu z=1$, which is all the contact-set argument needs. The formalisation proves the proposition, Corollary~\ref{cor:iso-smooth-components} and Theorem~\ref{thm:sharp-isoperimetric} in this $C^3$ form (the smooth case being a special case), and Step~2 of Theorem~\ref{thm:isoperimetric-rigidity} applies it to a $C^{3,\beta}$ representative.

\begin{lemma}[Strict superadditivity of $t^{2/3}$]\label{lem:concavity-23}
For $a,b>0$, $a^{2/3}+b^{2/3}>(a+b)^{2/3}$.
\end{lemma}

\begin{proof}
Strict concavity of $t\mapsto t^{2/3}$ on $(0,\infty)$ with value $0$ at $0$.
\end{proof}

\begin{corollary}[Smooth, possibly disconnected]\label{cor:iso-smooth-components}
\uses{prop:iso-smooth,lem:concavity-23}
For a bounded smooth domain $G$ with components $G_i$,
\[
  \Per(G)\ge(36\pi)^{1/3}\sum_i|G_i|^{2/3}\ge(36\pi)^{1/3}|G|^{2/3},
\]
and the second inequality is strict unless there is exactly one component.
\end{corollary}

\begin{proof}\uses{prop:iso-smooth,lem:concavity-23}
Apply \eqref{eq:iso-smooth} to each component and then
Lemma~\ref{lem:concavity-23}.
\end{proof}

\begin{theorem}[Sharp isoperimetric inequality]\label{thm:sharp-isoperimetric}
\uses{def:perimeter}
For every measurable $E\subset\R^3$ of finite measure and finite perimeter,
\begin{equation}\label{eq:sharp-isoperimetric}
  \Per(E)\ge c_I\,|E|^{2/3},
  \qquad
  c_I:=(36\pi)^{1/3}.
\end{equation}
\end{theorem}

\begin{proof}\uses{cor:iso-smooth-components,thm:smooth-approx}
Apply Corollary~\ref{cor:iso-smooth-components} to the smooth exact-volume
approximants of Theorem~\ref{thm:smooth-approx} and pass to the limit.
\end{proof}

\begin{remark}[The equality case is not proved here]\label{rem:iso-equality-later}
Equality in \eqref{eq:sharp-isoperimetric} holds exactly for balls up to null
sets; that is Theorem~\ref{thm:isoperimetric-rigidity}, proved in
Chapter~\ref{ch:rigidity} after the regularity theory.
\end{remark}

